% ============================================================
% Paper deep reading: DenseGRPO
% ============================================================

\section{论文精读：DenseGRPO}

\begin{conceptbox}
\textbf{论文：} \href{https://arxiv.org/abs/2601.20218}{DenseGRPO: From Sparse to Dense Reward for Flow Matching Model Alignment}\\
\textbf{版本：} arXiv 2601.20218v2，ICLR 2026。\\
\textbf{方向标签：} Flow Matching、Text-to-Image、GRPO、Dense Reward、Credit Assignment。\\
\textbf{阅读日期：} 2026-04-27。\\
\textbf{前置：} 需要已经理解 Flow-GRPO 的 SDE sampling、group relative advantage、policy ratio 和 KL regularization。
\end{conceptbox}

\subsection{0. 一句话结论}

DenseGRPO 解决的是 Flow-GRPO 里的核心信用分配问题：\textbf{最终图像 reward 不应该被原封不动分给所有 denoising steps}。它把终局 reward 改造成每个 timestep 的 dense reward，并进一步用 dense reward 来调每个 timestep 的探索噪声。

\begin{summarybox}
\textbf{我的判断：} 这是 Flow-GRPO 后最值得优先精读的一篇。因为它不是单纯堆一个技巧，而是把“reward 应该分给哪一步”这个问题具体化为可计算的 \key{step-wise reward gain}。
\end{summarybox}

\begin{figure}[htbp]
  \centering
  \includegraphics[width=\linewidth]{densegrpo/fig1_sparse_vs_dense_reward.png}
  \caption{Sparse reward 与 DenseGRPO dense reward 的核心差异。左侧是把 terminal reward 直接分给所有中间步骤；右侧是对每个 denoising step 估计独立贡献。图源：arXiv HTML。}
  \label{fig:densegrpo-sparse-vs-dense}
\end{figure}

\subsection{1. 这篇论文为什么重要}

Flow-GRPO 已经把 GRPO 接到了 flow matching / diffusion-like 视觉生成上。它的基本思路是：

\[
\text{prompt } c
\rightarrow
\text{sample a group of trajectories}
\rightarrow
\text{score final images}
\rightarrow
\text{GRPO update}
\]

问题在于，Flow-GRPO 的 reward 是 terminal reward：

\[
R(x_0, c)
\]

它只评价最终生成图像。但训练时每条轨迹有很多 denoising steps：

\[
x_T \rightarrow x_{T-1} \rightarrow \cdots \rightarrow x_1 \rightarrow x_0
\]

如果把同一个最终 reward 分给每一个 timestep，就隐含了一个很强的假设：

\[
\text{每一步对最终结果的贡献相同}
\]

这个假设通常是错的。某些步骤决定构图、颜色、对象关系；某些步骤只是细节 refinement。DenseGRPO 的核心价值，就是把这件事从直觉变成训练信号。

\subsection{2. Flow-GRPO 的问题在哪里}

Flow-GRPO 中，同一个 prompt 会采样一个 group。对第 \(i\) 个样本，最终图像是 \(x^i_0\)，reward model 给出：

\[
R(x^i_0, c)
\]

然后 group relative advantage 大致是：

\[
\hat{A}^i
=
\frac{
R(x^i_0,c)-\mathrm{mean}\{R(x^j_0,c)\}_{j=1}^{G}
}{
\mathrm{std}\{R(x^j_0,c)\}_{j=1}^{G}
}
\]

关键问题是：这个 \(\hat{A}^i\) 不依赖 timestep。也就是说，在第 \(t\) 步更新 policy 时，用的还是同一个整图 reward。

\begin{warnbox}
\textbf{问题本质：} Flow-GRPO 优化的是每一步 transition 的 logprob，但 advantage 来自整条 trajectory 的最终分数。这会导致 feedback-contribution mismatch。
\end{warnbox}

更具体地说，图像最后得分高，可能是因为早期 layout 选对了；也可能是因为后期文字细节修好了。把同一个 reward 分给所有步骤，会让模型不知道到底该强化哪一步。

\subsection{3. DenseGRPO 的核心想法}

DenseGRPO 要做两件事：

\begin{enumerate}[leftmargin=1.5em]
  \item \textbf{把 sparse terminal reward 变成 step-wise dense reward。}
  \item \textbf{根据 dense reward 分布，给不同 timestep 设置不同探索噪声。}
\end{enumerate}

这两件事对应论文的两个关键模块：

\[
\text{Step-wise Dense Reward}
\quad + \quad
\text{Exploration Space Calibration}
\]

\subsection{4. 模块一：Step-wise Dense Reward}

\subsubsection{4.1 从“最终分数”变成“每步增益”}

DenseGRPO 定义每一步的 reward 为相邻 latent 对应 reward 的提升量：

\[
\Delta R^i_t
=
R^i_{t-1} - R^i_t
\]

直觉是：

\[
x^i_t \rightarrow x^i_{t-1}
\]

这一小步 denoising 如果真的让图像更符合 reward model，那么从 \(x^i_t\) 对应的 clean image 到 \(x^i_{t-1}\) 对应的 clean image，reward 应该上升。这个上升量就是该步骤的贡献。

\begin{intubox}
\textbf{直觉版：} 不再问“最终图像好不好”，而是问“从这一步到下一步，图像有没有变得更好”。这就是 dense reward。
\end{intubox}

\subsubsection{4.2 中间 latent 怎么打分}

难点是：reward model 通常只能给 clean image 打分，不能直接给 noisy latent \(x_t\) 打分。

DenseGRPO 的办法是利用 flow matching ODE sampler 的确定性。给定中间 latent \(x^i_t\)，如果从它继续走 deterministic ODE denoising，可以得到一个 clean latent：

\[
\hat{x}^i_{t,0}
=
\mathrm{ODE}_n(x^i_t, c)
\]

再把 \(\hat{x}^i_{t,0}\) decode 成 clean image，并用普通 reward model 打分：

\[
R^i_t
:=
R(\hat{x}^i_{t,0}, c)
\]

这样就把“中间 latent 的 reward”转化成了“从这个 latent 确定性还原出来的 clean image 的 reward”。

\begin{summarybox}
\textbf{关键点：} DenseGRPO 没有额外训练 process reward model，也没有训练 critic。它复用已有 reward model，只是通过 ODE denoising 把中间 latent 映射回 clean image 后再打分。
\end{summarybox}

\subsubsection{4.2.1 疑问澄清：ODE denoising 不是探索}

这里容易产生一个误解：如果对某个 timestep 的 latent \(x_t\) 采取 ODE denoising，那 ODE 是 deterministic 的，怎么探索？

答案是：\textbf{ODE denoising 这一步本身不是用来探索的}。DenseGRPO 里有两条路径，需要分开看。

\textbf{第一条是采样/探索路径：SDE trajectory。}

训练时真正产生 rollout 的是 SDE sampler：

\[
x_T \rightarrow x_{T-1} \rightarrow \cdots \rightarrow x_0
\]

这里每一步有随机噪声，所以可以探索。Group 里不同样本、不同 trajectory 的差异，来自 SDE 注入噪声。

\textbf{第二条是打分/归因路径：ODE readout。}

对于已经由 SDE 采样得到的某个中间 latent \(x_t\)，DenseGRPO 额外做一次 deterministic ODE denoising：

\[
x_t
\xrightarrow{\mathrm{ODE}}
\hat{x}_{t,0}
\]

这一步不是为了探索，而是为了回答一个评估问题：

\begin{intubox}
如果当前已经走到这个 \(x_t\)，并且之后不再随机，只按确定性路径生成，它大概会通向什么 clean image？
\end{intubox}

然后对这个 clean prediction 打 reward：

\[
R_t = R(\hat{x}_{t,0}, c)
\]

所以 ODE 的角色是 \key{probe / readout / value estimation}，不是 policy rollout。

\[
\text{SDE = 产生探索样本}
\qquad
\text{ODE = 把中间 latent 映射成可打分的 clean prediction}
\]

DenseGRPO 的 dense reward：

\[
\Delta R_t = R_{t-1} - R_t
\]

可以理解为：沿着真实 SDE trajectory，从 \(x_t\) 走到 \(x_{t-1}\) 后，如果分别用 ODE 读出它们未来会变成什么 clean image，那么 clean-image reward 提升了多少。

\begin{warnbox}
\textbf{不要把两件事混在一起：} SDE 负责探索，ODE 负责评估。DenseGRPO 不是用 ODE 探索，而是先用 SDE 得到中间 latent，再用 ODE 把这些 latent 变成 reward model 可以评价的 clean-image proxy。
\end{warnbox}

\begin{figure}[htbp]
  \centering
  \includegraphics[width=\linewidth]{densegrpo/fig2_method_overview.png}
  \caption{DenseGRPO 总体流程：对中间 latent 做 ODE denoising，得到对应 clean image 的 reward，再用相邻 reward gain 构造 dense reward。图源：arXiv HTML。}
  \label{fig:densegrpo-method-overview}
\end{figure}

\subsubsection{4.3 新 advantage 怎么算}

拿到每个样本、每个 timestep 的 dense reward gain \(\Delta R^i_t\) 后，DenseGRPO 不再用整图 reward 做 group normalization，而是在每个 timestep 上单独做 group normalization：

\[
\hat{A}^i_t
=
\frac{
\Delta R^i_t - \mathrm{mean}\{\Delta R^j_t\}_{j=1}^{G}
}{
\mathrm{std}\{\Delta R^j_t\}_{j=1}^{G}
}
\]

这一步非常重要。它意味着：

\[
\text{第 } t \text{ 步的 policy update}
\quad
\text{只看第 } t \text{ 步的相对贡献}
\]

于是，不同 timestep 可以得到不同方向、不同强度的学习信号。

\subsection{5. 模块二：Exploration Space Calibration}

Flow-GRPO 为了让 deterministic flow model 可以做 RL 探索，把 ODE sampler 改成 SDE sampler，在每一步注入噪声：

\[
\sigma_t \epsilon
\]

原来的做法通常用一个统一噪声强度 \(a\)。DenseGRPO 认为这也不合理：不同 timestep 的作用不同，所需探索强度也不同。

\subsubsection{5.1 为什么统一噪声不对}

如果某个 timestep 噪声太大，生成轨迹会偏离合理分布，dense reward 可能几乎全是负的；如果噪声太小，group 内样本差异不够，RL 没有有效探索。

论文用 dense reward 的正负分布来判断探索是否合适：

\[
\Delta R^i_t > 0
\quad \text{vs.} \quad
\Delta R^i_t < 0
\]

如果一个 timestep 上正负样本比较均衡，说明该步既有好探索也有坏探索，训练信号更有区分度；如果几乎全正或全负，说明探索空间不合适。

\subsubsection{5.2 校准规则}

DenseGRPO 为每个 timestep 学一个噪声强度函数 \(\psi(t)\)。校准逻辑很简单：

\begin{itemize}[leftmargin=1.5em]
  \item 如果某个 timestep 的 positive / negative dense rewards 比较均衡，就轻微增加 \(\psi(t)\)，鼓励更多探索。
  \item 如果不均衡，就轻微降低 \(\psi(t)\)，避免过度或无效探索。
\end{itemize}

最终采样时使用：

\[
\sigma_t = \psi(t)
\]

而不是所有 timestep 共享一个标量 \(a\)。

\begin{intubox}
\textbf{直觉版：} Dense reward 不只是训练 signal，它还变成了诊断工具：用来判断每个时间步的探索噪声是否合适。
\end{intubox}

\begin{figure}[htbp]
  \centering
  \includegraphics[width=\linewidth]{densegrpo/fig3_dense_reward_noise_calibration.png}
  \caption{不同噪声强度下的 dense reward 分布。固定噪声会在部分 timestep 造成奖励极度不平衡；DenseGRPO 用 timestep-specific \(\psi(t)\) 校准探索空间。图源：arXiv HTML。}
  \label{fig:densegrpo-noise-calibration}
\end{figure}

\subsection{6. 整体训练流程}

可以把 DenseGRPO 的训练理解成下面几步：

\begin{enumerate}[leftmargin=1.5em]
  \item 对同一个 prompt 采样 \(G\) 条 SDE trajectories。
  \item 对每条 trajectory 的每个 intermediate latent \(x^i_t\)，用 ODE denoising 得到 \(\hat{x}^i_{t,0}\)。
  \item 对 \(\hat{x}^i_{t,0}\) decode 后用 reward model 打分，得到 \(R^i_t\)。
  \item 用相邻 reward 差值计算 step-wise dense reward：
  \[
  \Delta R^i_t = R^i_{t-1} - R^i_t
  \]
  \item 对每个 timestep，在 group 内 normalize \(\Delta R^i_t\)，得到 \(\hat{A}^i_t\)。
  \item 用原来的 GRPO clipped objective 更新 policy，只是 advantage 从 trajectory-level 换成 step-level。
  \item 训练前用 dense reward 分布校准 \(\psi(t)\)，让每个 timestep 的 SDE exploration 更合理。
\end{enumerate}

\subsection{7. 实验看什么}

论文沿用 Flow-GRPO 的三类 text-to-image 任务：

\begin{itemize}[leftmargin=1.5em]
  \item \textbf{Compositional Image Generation：} reward model 用 GenEval。
  \item \textbf{Human Preference Alignment：} reward model 用 PickScore。
  \item \textbf{Visual Text Rendering：} reward 来自 OCR accuracy。
\end{itemize}

主要对比：

\begin{itemize}[leftmargin=1.5em]
  \item Flow-GRPO：terminal reward 直接分配给所有 timestep。
  \item Flow-GRPO+CoCA：用 latent similarity 给 step-wise feedback。
  \item DenseGRPO：用 ODE-based intermediate clean image 估计 dense reward gain。
\end{itemize}

论文报告 DenseGRPO 在三类任务上都优于这些 baseline。比较关键的是，DenseGRPO 相对 Flow-GRPO+CoCA 也有明显收益，这说明它的优势不只是“用了 step-wise reward”，而是 \key{ODE-based latent reward 比 latent similarity 更贴近真实贡献}。

\begin{figure}[htbp]
  \centering
  \includegraphics[width=\linewidth]{densegrpo/fig4_main_results.png}
  \caption{主实验结果柱状图。DenseGRPO 在组合图像生成、人类偏好对齐和视觉文字渲染任务上都优于主要 baseline。图源：arXiv HTML。}
  \label{fig:densegrpo-main-results}
\end{figure}

\begin{figure}[htbp]
  \centering
  \includegraphics[width=0.92\linewidth]{densegrpo/fig5_qualitative_results.png}
  \caption{质性结果对比。论文用这些样例说明 DenseGRPO 在组合关系、语义一致性和文本渲染上有更好的对齐效果。图源：arXiv HTML。}
  \label{fig:densegrpo-qualitative}
\end{figure}

\subsection{8. Ablation 该怎么读}

\begin{center}
  \includegraphics[width=\linewidth]{densegrpo/fig6_ablation_analysis.png}

  {\small \textbf{图 6：} 消融实验。Dense reward、timestep-specific noise 和 ODE denoising steps 分别对应这里的三个核心问题：step-wise reward 是否有用、\(\psi(t)\) 是否比固定噪声更好、ODE readout 步数是否影响 reward 估计质量。图源：arXiv HTML。}
\end{center}

\subsubsection{8.1 Dense reward 是否真的有用}

论文做了一个 Dense Reward Baseline：不用 reward gain，只直接用 \(R^i_{t-1}\) 来优化第 \(t\) 步。它比 Flow-GRPO 好，说明 dense reward 本身有价值；但 DenseGRPO 更好，说明“相邻 reward gain”比“下一步 latent reward 绝对值”更合理。

\[
R^i_{t-1}
\quad \text{不如} \quad
R^i_{t-1} - R^i_t
\]

原因是绝对 reward 仍然混入了未来步骤的累计贡献，而 reward gain 更接近当前一步的边际贡献。

\subsubsection{8.2 Timestep-specific noise 是否有用}

论文比较了固定噪声 \(a=0.7\) 和校准后的 \(\psi(t)\)。结论是 \(\psi(t)\) 更好。

这说明 DenseGRPO 的第二个模块不是附属技巧：一旦你有了 step-wise dense rewards，就能看到哪些时间步探索过强或过弱，并据此调 sampling noise。

\subsubsection{8.3 ODE denoising steps 的成本与收益}

DenseGRPO 需要对每个 intermediate latent 做 ODE denoising。论文比较 \(n=1,2,t\)，发现 \(n=t\) 最好，但计算也更贵。

\[
\text{full readout cost}
\approx T + (T-1) + \cdots + 1
= \frac{T(T+1)}{2}
\]

也就是说，如果 rollout 本身有 \(T\) 个 denoising steps，而每个 \(x_t\) 都用完整 ODE 读出到 clean image，那么 dense reward 的额外模型调用会从线性的 \(O(T)\) 接近变成平方级 \(O(T^2)\)。这还没有算 VAE decode 和 reward model scoring 的额外开销。

论文在 Fig. 6(c) 里确实报告了训练成本：在相同实验设置下，\(n=1\)、\(n=2\)、\(n=t\) 分别需要 11、13、19 GPU hours 来训练 20 steps。也就是说，最准确的 \(n=t\) readout 明显更贵；论文的说法是更大的 \(n\) 虽然有更高计算开销，但同等 GPU training time 下收益更好。

\begin{warnbox}
\textbf{工程代价：} DenseGRPO 的核心收益来自更准的 dense reward，但代价是多次 ODE denoising、VAE decode 和 reward model scoring。它不是免费的加速方法，而是更精细、更贵的训练信号。对于视频模型，这个问题会被进一步放大，因为单次 denoising forward 和 reward evaluation 都远贵于图像模型。
\end{warnbox}

\textbf{我的判断：} full \(n=t\) 更像是一个证明 dense reward gain 上限有效的版本。实际做大图像模型或视频模型后训练时，更现实的实现可能是：只对少量关键 timesteps 做 readout，或者固定 \(n=1,2,k\) 做 truncated readout，再用插值、chunk 或 learned proxy 来近似 step-wise advantage。

\subsection{9. 和 TempFlow-GRPO / Chunk-GRPO 的关系}

这三篇都在回答 Flow-GRPO 的 credit assignment 问题，但切入点不同：

\begin{itemize}[leftmargin=1.5em]
  \item \textbf{TempFlow-GRPO：} 问“哪个 timestep 更关键”，用 trajectory branching 看 reward 差异。
  \item \textbf{Chunk-GRPO：} 问“step 太细、trajectory 太粗怎么办”，把连续 steps 合成 chunk。
  \item \textbf{DenseGRPO：} 问“每一步真实贡献怎么直接估计”，用 ODE-based reward gain 做 dense reward。
\end{itemize}

如果只读一篇来理解这条线，DenseGRPO 最适合作为核心，因为它最直接地把 terminal reward 拆成了 step-wise signal。

\subsection{10. 局限与风险}

\begin{itemize}[leftmargin=1.5em]
  \item \textbf{计算成本高：} 每个 timestep 的 intermediate latent 都可能需要 ODE denoising、VAE decode 和 reward scoring。若使用完整 \(n=t\) readout，额外 denoising 调用接近 \(T(T+1)/2\)，比普通 \(O(T)\) rollout 贵很多。
  \item \textbf{论文只给了局部成本：} 实验使用 16 张 NVIDIA A100。论文报告 \(n=1,2,t\) 训练 20 steps 分别需要 11、13、19 GPU hours，但没有给出完整 4500/1500/4500 iteration 训练的 wall-clock time。
  \item \textbf{依赖 reward model 质量：} 如果 reward model 有偏，dense reward 会把偏差更精细地传播到每一步。
  \item \textbf{可能放大奖励黑客：} 论文附录也提到，更准确的 step-wise reward 可能让模型更容易过拟合某些 reward model。
  \item \textbf{适用前提：} 需要能从 intermediate latent 通过 deterministic sampler 映射到 clean image。对没有稳定 deterministic path 的模型，需要额外改造。
\end{itemize}

\subsection{11. 对图像编辑后训练的启发}

DenseGRPO 对图像编辑很有启发，因为图像编辑的 reward 更不应该只看最终图：

\[
\text{edit region 应该变}
\quad
\text{preserve region 不应该变}
\]

如果未来做 image editing RL，可以借鉴 DenseGRPO 的想法：

\begin{itemize}[leftmargin=1.5em]
  \item 对每个 timestep 还原 intermediate clean image。
  \item 分别计算 edit reward、preservation reward、identity reward、aesthetic reward。
  \item 用相邻 timestep 的 reward gain 判断哪一步真正造成了目标区域变化或非目标区域破坏。
  \item 再把不同区域、不同 reward component 的 gain 分给对应 timestep。
\end{itemize}

\begin{summarybox}
\textbf{可复用点：} DenseGRPO 的真正可迁移贡献不是某个公式，而是一个范式：
\[
\begin{aligned}
\text{intermediate latent}
&\rightarrow \text{deterministic clean prediction}\\
&\rightarrow \text{reward model scoring}\\
&\rightarrow \text{step-wise reward gain}
\end{aligned}
\]
这个范式可以迁移到图像编辑、视频生成、甚至 discrete diffusion 的后训练。
\end{summarybox}

\subsection{12. 读完要记住的三个问题}

\begin{enumerate}[leftmargin=1.5em]
  \item \textbf{为什么 terminal reward 不能直接分给所有 timestep？}\\
  因为最终结果是所有 steps 的累计贡献，单步更新需要的是当前 step 的边际贡献。
  \item \textbf{DenseGRPO 怎么给 intermediate latent 打分？}\\
  用 deterministic ODE denoising 把 \(x_t\) 映射到 clean image，再用已有 reward model 打分。
  \item \textbf{为什么 dense reward 还能调探索噪声？}\\
  因为 dense reward 的正负分布可以反映某个 timestep 的探索是否过强或过弱。
\end{enumerate}
